
Our experiments demonstrate that an oracle gap of 15.09\% relative improvement (11.62 percentage points absolute) exists between per-task optimal adapter rank selection and the best fixed-rank baseline, with oracle selections distributed across all tested ranks.

\subsubsection{Main Results: Oracle Gap Measurement}

Finding: Oracle gap exceeds 10\% threshold, validating existence of task heterogeneity.

Table 1 presents oracle gap measurements across all 17 tasks.

\textbf{Table 1: Oracle Gap Measurement}

\begin{table}[htbp]
\centering
\begin{tabular}{lll}
\toprule
Configuration & Average Accuracy & Tasks Optimal \\
\midrule
Oracle (per-task best) & 88.58\% & — \\
Fixed rank-8 (best) & 76.97\% & 4 tasks \\
Fixed rank-16 & 75.37\% & 4 tasks \\
Fixed rank-4 & 73.66\% & 5 tasks \\
Fixed rank-32 & 62.95\% & 4 tasks \\
Oracle Gap (abs) & 11.62 pp & — \\
Oracle Gap (rel) & 15.09\% & — \\
\bottomrule
\end{tabular}
\end{table}

Note: Relative gap computed as (Oracle_avg - Best_fixed) / Best_fixed × 100\%. The 15.09\% represents a relative improvement over the baseline, equivalent to 11.62 percentage points in absolute terms.

Analysis:

\begin{enumerate}
\item Oracle gap (15.09\% relative improvement, 11.62 percentage points absolute) exceeds 10\% threshold by 5.09 percentage points. This validates the hypothesis that multi-domain benchmarks exhibit task heterogeneity creating measurable optimization opportunities. The gap represents an upper bound for task-aware routing—practical systems with imperfect classifiers will achieve a fraction of this improvement.
\end{enumerate}

\begin{enumerate}
\item Best fixed rank-8 achieves 76.97\% average accuracy. This aligns with literature consensus that rank 8 is a reasonable default choice. However, the oracle (88.58\%) significantly outperforms this configuration, demonstrating that no single rank serves all tasks optimally.
\end{enumerate}

\begin{enumerate}
\item Rank-32 performs worst (62.95\%) despite having highest capacity. The 14 percentage point gap between rank-32 and rank-8 demonstrates severe overfitting.
\end{enumerate}

\subsubsection{Heterogeneity Analysis: Oracle Selection Distribution}

Finding: Oracle selections distribute across ranks, indicating genuine task diversity.

If one rank dominated, most tasks would select it as oracle—indicating that a fixed configuration works adequately. Table 2 shows distribution across ranks.

\textbf{Table 2: Oracle Selection Distribution}

\begin{table}[htbp]
\centering
\begin{tabular}{lll}
\toprule
Rank & Tasks Selecting as Oracle & Example Tasks \\
\midrule
4 & 5 tasks (29\%) & CoLA, STS-B, WNLI, XNLI-zh, PAWS-X-zh \\
8 & 4 tasks (24\%) & SST-2, MNLI, XNLI-en, PAWS-X-en \\
16 & 4 tasks (24\%) & MRPC, QNLI, XNLI-es, PAWS-X-es \\
32 & 4 tasks (24\%) & QQP, RTE, XNLI-de, PAWS-X-de \\
\bottomrule
\end{tabular}
\end{table}

Analysis:

\begin{enumerate}
\item Distribution (5/4/4/4) approximates uniformity, providing evidence that different tasks prefer different capacity levels. No single rank dominates.
\end{enumerate}

\begin{enumerate}
\item Language-specific patterns emerge. Chinese cross-lingual tasks (XNLI-zh, PAWS-X-zh) both select rank-4, while German tasks (XNLI-de, PAWS-X-de) both select rank-32. This systematic structure suggests task meta-features correlate with optimal rank.
\end{enumerate}

\begin{enumerate}
\item Dataset size correlates with optimal rank. Small datasets (CoLA: 8.5K samples, WNLI: 635 samples) prefer low ranks (rank-4), while large datasets (QQP: 363K samples, MNLI: 392K samples) can leverage higher ranks without overfitting.
\end{enumerate}

\subsubsection{Comparative Analysis: Fixed-Rank Performance}

Finding: Fixed-rank performance varies substantially, with rank-32 collapsing on small and medium datasets.

Rank-32's poor average performance (62.95\%) warrants investigation. Table 3 breaks down rank-32 performance by dataset size.

\textbf{Table 3: Rank-32 Performance by Dataset Size}

\begin{table}[htbp]
\centering
\begin{tabular}{lllll}
\toprule
Dataset Size & Task & Rank-32 Accuracy & Rank-4 Accuracy & Gap \\
\midrule
Small (<10K) & CoLA (8.5K) & 50.00\% (random) & 86.88\% & -36.88 pp \\
Small (<10K) & WNLI (635) & 50.00\% (random) & 88.42\% & -38.42 pp \\
Medium (10K-100K) & SST-2 (67K) & 50.00\% (random) & 81.20\% & -31.20 pp \\
Large (>100K) & QQP (363K) & 88.49\% & 50.00\% & +38.49 pp \\
Large (>100K) & MNLI (392K) & 55.22\% & 86.05\% & -30.83 pp \\
\bottomrule
\end{tabular}
\end{table}

Analysis:

\begin{enumerate}
\item Rank-32 collapses to random baseline (50\%) on small and medium datasets. The collapse occurs not just on tiny datasets (WNLI: 635 samples, CoLA: 8.5K samples) but also on medium-sized datasets like SST-2 with 67K samples. This demonstrates that rank-32 requires substantially larger datasets to avoid catastrophic overfitting.
\end{enumerate}

\begin{enumerate}
\item Rank-32 achieves high performance only on the largest dataset (QQP: 363K samples), achieving 88.49\% and substantially outperforming rank-4's 50.00\%. However, even MNLI with 392K samples shows rank-32 underperforming rank-4 (55.22\% vs 86.05\%), suggesting that dataset size alone doesn't determine optimal rank.
\end{enumerate}

\begin{enumerate}
\item Rank-32 exhibits extreme brittleness. Performance varies from 50\% (random baseline) to 88.49\% across tasks, demonstrating high variance and unpredictable behavior. In contrast, ranks 4-16 show more consistent performance across dataset sizes.
\end{enumerate}

\begin{enumerate}
\item The overfitting threshold is higher than expected. Previous work suggested rank-32 might work for datasets >10K samples with proper regularization. Our results show that even 67K samples (SST-2) are insufficient—rank-32 collapses to random baseline.
\end{enumerate}
